\chapter{INTRODUCTION}
\label{ch:introduction}

\section{Introduction}

Electrospinning is a fabrication technique in which a high electric field applied
between a spinneret and a grounded collector draws a polymer solution into fibres
with diameters ranging from tens of nanometres to a few micrometres. The non-woven
mats obtained in this way combine a very high surface-to-volume ratio with an
interconnected pore structure, and these properties explain their use in air
and water filtration, in wound dressings, in drug delivery and as scaffolds for
tissue engineering.

The process, however, is not uniformly well behaved. When the electrostatic
stretching force is not sufficient to overcome the surface tension of the jet, the
solution collapses locally into droplets which then solidify as ellipsoidal
swellings along the fibre. These \emph{bead defects}, first characterised
systematically by Fong, Chun and Reneker \cite{fong1999beaded}, are commonly taken
as an indicator that the spinning parameters lie outside their usable window. Their
presence reduces the effective surface area of the mat, disturbs its pore-size
distribution and introduces stress concentrators that weaken it mechanically. For
this reason, counting and measuring the beads is a routine step whenever an
electrospinning recipe is being tuned.

That step is currently performed manually. The operator loads each scanning
electron microscopy (SEM) micrograph into ImageJ \cite{schindelin2012fiji},
identifies the beads by visual inspection and outlines them one at a time. A single
micrograph of the dataset used here contains up to 152 beads, each of which has to
be traced individually. The procedure is therefore slow and it does not scale to
the number of micrographs that a parameter study produces. In addition, the
boundary between a large bead and a locally thick fibre segment is a matter of
judgement, so that the resulting annotation is difficult to reproduce between
operators or even between sessions of the same operator.

\section{Background and Motivation}

The automation of this measurement is an image segmentation problem, although not
among the simpler ones. Three properties of SEM micrographs of nanofibre mats make
the task more difficult than the generic case: i) a dense and high-contrast
background, ii) objects of interest that touch and overlap, and iii) scarce data
acquired at several magnifications. These are considered in turn below.

The beads are located on top of a dense background of overlapping fibres. The
fibres are themselves bright, elongated and cross each other at every angle, and
consequently a simple intensity threshold is not able to separate beads from
background, as both occupy similar grey levels while the texture energy of the
fibre mesh is comparable to that of a bead surface.

The objects of interest touch and overlap. Beads frequently occur in clusters along
the same fibre, and a method that produces only a foreground mask merges such a
cluster into a single connected component. Since the quantities required by the
laboratory are a \emph{count} and a \emph{size distribution}, the merging of two
beads into one cannot be treated as a small error, because it lowers the reported
count and at the same time introduces into the size histogram a large particle that
does not exist. It is this requirement that makes the task an instance segmentation
problem rather than a semantic one.

The data are scarce and are acquired at several magnifications. The dataset
available for this thesis contains 11 annotated micrographs drawn from four
specimens, each of them imaged at $500\times$, $1000\times$ and $3000\times$. A
bead measuring approximately 14 pixels across at $500\times$ measures
approximately 76 pixels across at $3000\times$, so that the same physical object
appears with a sixfold difference in apparent size. A model trained on such data is
required to be scale-robust, while the evaluation protocol has to take into account
that the four specimens, and not the eleven images, are the independent units of
the experiment.

The first two difficulties are addressed directly by modern instance segmentation
architectures. Mask R-CNN \cite{he2017mask} detects each object separately and
predicts a mask inside its own region proposal, so that touching beads remain
separate instances by construction. The third difficulty, namely the limited amount
of data and the manner in which it is evaluated, determines whether the reported
numbers can be trusted, and for this reason it is given explicit attention
throughout this thesis.

\section{Problem Statement}

This thesis addresses the automated detection, delineation and quantification of
bead defects in SEM micrographs of electrospun nanofibre mats. The work is
organised around four research questions:

\begin{enumerate}
    \item Is a deep instance segmentation model able to delineate individual bead
    defects in SEM micrographs of nanofibre mats with an accuracy sufficient to
    replace manual annotation, when only 11 annotated micrographs from four
    specimens are available?

    \item Does instance-level modelling offer a measurable advantage over semantic
    segmentation in this task, in particular in the separation of touching beads
    and consequently in the accuracy of the reported bead count?

    \item How does a deep model compare with the classical image-processing
    pipeline, consisting of thresholding followed by watershed splitting, which
    represents current laboratory practice?

    \item Do the derived physical quantities, namely bead count, areal density and
    size distribution expressed in micrometres, agree with those obtained from the
    manual annotation closely enough to be used in process optimisation?
\end{enumerate}

The scope of the work is restricted to bead defects in micrographs of a single
material system, imaged on one instrument. Fibre diameter measurement,
three-dimensional structure and the relation between spinning parameters and bead
formation are outside this scope, although they are discussed as possible
extensions in Section~\ref{sec:futurework}.

\section{Research Objectives}
\label{sec:objectives}

The objectives are listed below, each together with the criterion by which it
is judged to have been met.

\begin{itemize}
    \item \textbf{Dataset characterisation and calibration.} Establish what the
    available data contain and place every measurement on a physical scale. This
    includes the recovery of the pixel size of each micrograph from its burned-in
    scale bar and the verification of this calibration by checking whether the bead
    size distribution is consistent across the three magnifications. Success
    criterion: median bead diameter agreeing to within 20\% across $500\times$,
    $1000\times$ and $3000\times$.

    \item \textbf{A leakage-free evaluation protocol.} Define a partitioning scheme
    appropriate to a dataset whose 11 images derive from only 4 specimens, and
    demonstrate quantitatively why the naive alternative, that is a random split
    over images or over the augmented copies, overstates performance. Success
    criterion: a reported gap between the naive and the specimen-wise protocol
    obtained with the same model.

    \item \textbf{Instance segmentation of beads.} Train Mask R-CNN to delineate
    individual beads under leave-one-specimen-out cross-validation. Success
    criterion: $\mathrm{AP}_{50} \geq 0.30$ averaged over the four folds, which
    corresponds to what a standard region-based detector achieves on a published
    benchmark of comparable object size.

        This criterion replaces the value of $\mathrm{AP}_{50} \geq 0.70$ that was
    recorded when the work was proposed. The grounds for the revision are reported
    here explicitly, since a criterion that is revised after the results are known
    has no value unless the grounds for revising it are independent of those
    results. The original figure was set before the small-object literature
    reviewed in Section~\ref{sec:small_objects} was surveyed, and that survey
    indicates that 0.70 lies above what any published detector reaches at this
    object size. On AI-TOD-v2, a benchmark whose objects average 12.7 pixels
    against the 14.3-pixel median obtained here, a standard Faster R-CNN reaches
    $\mathrm{AP}_{50} = 0.299$, whereas the method introduced in the same paper
    reaches 0.574 \cite{xu2022tiny}. The revised criterion is the first of these
    values, rounded. It is stated while being fully aware that
    Chapter~\ref{ch:results} reports 0.298, which is two thousandths below it, as a
    threshold selected in order to be cleared would not have been set at this
    level.

    One qualification has to be attached to this number, and it would soften the
    result. The AI-TOD-v2 figure is a \emph{box} average precision, whereas the
    value of 0.298 reported here is a \emph{mask} average precision, which is the
    stricter of the two since a mask is required to agree in shape as well as in
    position. When measured on boxes, as is done in the reference, the
    detection-only methods of Chapter~\ref{ch:results} do clear the criterion,
    reaching 0.305 and 0.315. The criterion is nevertheless kept at the mask level,
    where it is not met, because the objective stated here is instance segmentation
    and not detection.

    \item \textbf{Comparison against five alternatives.} Quantify what is gained by
    instance-level modelling and which component of it is responsible for that
    gain. Mask R-CNN is compared against i) the same detector with its mask branch
    removed, that is Faster R-CNN \cite{ren2015faster}, which isolates the
    contribution of that branch, ii) two one-stage detectors of the YOLO family
    \cite{jocher2023ultralytics}, a YOLOv8 segmentation model and a YOLOv5u
    detector, which vary the architecture family as a whole and differ from each
    other in the backbone, iii) a U-Net \cite{ronneberger2015unet} semantic
    baseline followed by connected-component labelling, and iv) an Otsu
    \cite{otsu1979threshold} and watershed \cite{vincent1991watersheds} pipeline
    that represents current manual practice. Success criterion: a clear ranking on
    counting error, with the merge behaviour of each method reported explicitly.

    \item \textbf{Physical quantification.} Report bead count, areal density
    (beads\,mm$^{-2}$) and the size distribution in micrometres, and compare each
    of them against the manual ground truth. Success criterion: counting error
    together with a distributional comparison of predicted versus manual bead
    diameters.
\end{itemize}

\section{Contributions}

\begin{enumerate}
    \item A characterised and physically calibrated dataset of 606 manually
    delineated bead defects across 11 SEM micrographs of four electrospun
    specimens, in which the per-magnification pixel sizes are recovered from the
    instrument's own scale bars and are validated for internal consistency.

    \item An evaluation protocol for this class of small and hierarchically
    structured microscopy datasets, in which the specimen and not the image is used
    as the unit of partitioning, together with a measurement of the cost that a
    naive split actually carries on a dataset of this shape.

    \item A controlled comparison of six methods, namely two-stage instance
    segmentation, the same detector without its mask branch, two one-stage
    alternatives from the same family, semantic segmentation with connected
    components and classical thresholding, carried out on the same micrographs and
    under the same protocol. The comparison is evaluated both by segmentation
    metrics and by the counting and sizing errors that determine whether a method
    is usable in the laboratory.

    \item An open implementation covering annotation conversion, scale calibration,
    training, evaluation and the reporting of physical quantities.
\end{enumerate}

\section{Thesis Organization}

\textbf{Chapter 2} reviews electrospinning and bead formation, classical image
analysis for fibrous microscopy and deep segmentation architectures, and identifies
the gap that this thesis addresses.

\textbf{Chapter 3} presents the methodology, that is the dataset and its structure,
the scale calibration, the specimen-wise evaluation protocol, the six methods under
comparison and the metrics.

\textbf{Chapter 4} describes the implementation, including annotation conversion,
the handling of the burned-in information banner, the augmentation strategy and the
training setup.

\textbf{Chapter 5} reports the experimental results and discusses them.

\textbf{Chapter 6} concludes and sets out directions for future work.
